\scanpage{170}
\mspage{6}
% archive sheet number: 169
2. \uline{Réduction à un cas \unsure{standard|standart} de dim. relative 1}.

\deleted{La réduction de 1.3. Au cas où $X$ est lisse sur $K$, de dim relative 1 se fait suivant les lignes de la réduction analogue faite dans II 3, du moins}

2.1. \uline{Réduction de 1.3 au cas} $X$ \deleted{\unsure{affine géom.}} \deleted{\unsure{géométriquement projectif}} \deleted{\uline{géom. et normal}} \inserted{\deleted{\unsure{affine}} \uline{géométriquement normal et géom intègre}}

\deleted{On recouvre $X$ par des ouverts affines $X_i$}
\deleted{2.1.1. Réduction au cas $X$ affine}
\deleted{en nb fini, on pose $X'=\coprod_i X_i$, $X''=X'\mathbin{\times}_X X'$, $X'''=X'\mathbin{\times}_X X'\mathbin{\times}_X X'$.}

\deleted{\unsure{Soit} $X$} Notons que l'on peut \deleted{\unsure{supposer provisoirement}} \inserted{supposer $S$} complet (I\quad), et qu'il est inoffensif pour la conclusion de remplacer $V$ par \unsure{sa|son} normalisé dans une extension finie $K'$ de $K$, et $X_K$ par $X_{K'}=X_K\otimes_K K'$ (car cela revient à remplacer $I$ par un s-groupe ouvert $I'$).
On peut alors, quitte à passer à une extension radicielle $K'$ de $K$, supposer que \deleted{\unsure{$X$}} le normalisé \deleted{\unsure{de}} de $(X_K)_{\mathrm{red}}$ est \uline{géométriquement normal} (EGA \uline{IV} \unsure{17.11.14.2|17, II.14.2}).
\deleted{De plus (EGA IV 4) on peut supposer $X_K$ géom. connexe. Enfin}

% slanted
\margin{\unsure{on est ramené} au cas où le pt de base sur $X_K$ est un $x\in X_K(K)$}
Soit $X_0$ le normalisé de $X_{\mathrm{red}}$ et soient $X_1$ et $X_2$ ses produits fibrés \inserted{sur $X$}. Dans SGA 1 IX 5, nous avons donné une méthode à la Van Kampen pour calculer le groupe fondamental de $X$ par générateurs et relations en termes de ceux des \inserted{comp. connexes des} $X_0,X_1,X_2$. Nous pouvons d'ailleurs supposer, quitte à faire une ext. finie sur $K$, que ces comp. connexes sont géom. connexes (EGA IV 4). \deleted{Alors le groupe fondamental} \inserted{\unsure{Nous pouvons aussi}} supposer que les pts bases choisis dans ces comp. \unsure{supposées|sont}
